%!TEX program = xelatex
\documentclass[lang=en,11pt,a4paper,citestyle=authoryear]{elegantpaper}

\title{Homework05 - MATH 720}
\author{Boren(Wells) Guan}
\date{April 21, 2026}

\usepackage{array,url}
\usepackage[notextcomp]{stix}
\usepackage{subfigure}
\usepackage{mathtools,amsmath,amssymb,amsthm,mathrsfs}
\usepackage{bbm}
\newcommand{\ccr}[1]{\makecell{{\color{#1}\rule{1cm}{1cm}}}}
\newcommand{\code}[1]{\lstinline{#1}}
\newcommand{\prvd}{$\hfill \qedsymbol$}
\newcommand{\Z}{\mathbb{Z}}
\newcommand{\R}{\mathbb{R}}
\newcommand{\N}{\mathbb{N}}
\newcommand{\C}{\mathbb{C}}
\newcommand{\Q}{\mathbb{Q}}
\newcommand{\M}{\mathcal{M}}
\newcommand{\B}{\mathcal{B}}
\newcommand{\X}{\mathcal{X}}
\newcommand{\Hil}{\mathcal{H}}
\newcommand{\range}{\mathcal{R}}
\newcommand{\nul}{\mathcal{N}}
\newcommand{\ip}[2]{\langle #1,#2\rangle}
\newcommand{\supp}{\operatorname{supp}}
\newcommand{\dist}{\operatorname{dist}}
\newcommand{\Om}{\Omega}
\newcommand{\pa}{\partial}
\newcommand{\Boxg}{\square_g}
\newcommand{\Boxop}{\square}
\newcommand{\sgn}{\operatorname{sgn}}
\newcommand{\Ree}{\operatorname{Re}}
\newcommand{\Ima}{\operatorname{Im}}
\newcommand{\1}{\mathbbm{1}}
\newcommand{\ov}{\overline}
\newcommand{\la}{\lambda}
\newcommand{\ds}{\displaystyle}

\begin{document}
\maketitle

\subsection*{Before Reading:}\par
As in the previous homework, I will focus on writing the core arguments in a reasonably clean way, and I will use standard facts from PDE whenever they are routine. The main goal here is to organize the multiplier and energy arguments so that each estimate is transparent.

\subsection*{Problem 1}
Consider the wave equation
\[
\Boxop u:=\partial_t^2u-\Delta_xu=0
\]
in a cylinder $\Omega\times[0,T]$, with initial data
\[
u(0)=u_0,\qquad \partial_tu(0)=u_1,
\]
and Dirichlet boundary condition
\[
u(t,x)=0\qquad \text{on }[0,T]\times\partial\Omega.
\]
Prove that energy estimates hold.\par
Extra challenge: Prove the analogous statement for
\[
g^{\alpha\beta}\partial_\alpha\partial_\beta u=0
\]
with Lipschitz metric under the assumptions that the slices $t=\mathrm{const}$ are spacelike and the lateral boundary $[0,T]\times\partial\Omega$ is timelike.

\vspace{0.5em}
\textbf{Sol.} \par
Define the energy
\[
E(t):=\frac12\int_\Omega \bigl(|\partial_tu(t,x)|^2+|\nabla_xu(t,x)|^2\bigr)\,dx.
\]
We claim that $E(t)$ is constant on $[0,T]$, i.e.
\[
E(t)=E(0)=\frac12\int_\Omega \bigl(|u_1|^2+|\nabla u_0|^2\bigr)\,dx.
\]
since we may multiply the equation $\partial_t^2u-\Delta u=0$ by $\partial_tu$ and then we have
\[
\int_\Omega \partial_t^2u\,\partial_tu\,dx-\int_\Omega \Delta u\,\partial_tu\,dx=0.
\]
The first term:
\[
\int_\Omega \partial_t^2u\,\partial_tu\,dx=\frac12\frac{d}{dt}\int_\Omega |\partial_tu|^2\,dx.
\]
For the second term, by the integration by parts in space, we have
\[
-\int_\Omega \Delta u\,\partial_tu\,dx
=\int_\Omega \nabla u\cdot \nabla\partial_tu\,dx-\int_{\partial\Omega}\partial_\nu u\,\partial_tu\,dS.
\]
By the Dirichlet condition, the boundary term vanishes and also
\[
\int_\Omega \nabla u\cdot \nabla\partial_tu\,dx
=\frac12\frac{d}{dt}\int_\Omega |\nabla u|^2\,dx.
\]
Then we have
\[
\frac{d}{dt}E(t)=0.
\]
Therefore, the energy is conserved, and in particular
\[
\|\partial_tu(t)\|_{L^2(\Omega)}+\|\nabla u(t)\|_{L^2(\Omega)}
\lesssim \|u_1\|_{L^2(\Omega)}+\|\nabla u_0\|_{L^2(\Omega)}
\]
for all $t\in[0,T]$, so the energy estimate goes.

For the variable coefficient problem, it is better to use the stress-energy tensor. Write
\[
\Boxg u:=g^{\alpha\beta}\partial_\alpha\partial_\beta u=0.
\]
Associated to $u$ we consider
\[
T_{\alpha\beta}[u]:=\partial_\alpha u\,\partial_\beta u-
\frac12 g_{\alpha\beta}g^{\mu\nu}\partial_\mu u\,\partial_\nu u.
\]
When $u$ solves $\Boxg u=0$, the standard computation gives
\[
\nabla^\alpha T_{\alpha\beta}=0.
\]
Now choose a future timelike vector field $X$ which is transversal to the hypersurfaces $t=\mathrm{const}$; the natural choice is $X=\partial_t$ in the given coordinate system. Define the current
\[
P_\alpha:=T_{\alpha\beta}X^\beta.
\]
Then
\[
\nabla^\alpha P_\alpha=T_{\alpha\beta}\nabla^\alpha X^\beta.
\]
Since $g$ is Lipschitz, $\nabla X\in L^\infty$, so the right-hand side is bounded by
\[
|\nabla^\alpha P_\alpha|\lesssim |\partial u|^2.
\]
Integrate this divergence identity in the slab
\[
D_t:=\{(s,x):0\le s\le t,\ x\in\Omega\}.
\]
By the divergence theorem, the boundary of $D_t$ consists of the initial slice $\Sigma_0$, the final slice $\Sigma_t$, and the lateral boundary $\Gamma=[0,t]\times\partial\Omega$. Thus
\[
\int_{\Sigma_t} T_{\alpha\beta}n^\alpha X^\beta\,d\sigma
-\int_{\Sigma_0} T_{\alpha\beta}n^\alpha X^\beta\,d\sigma
+\int_{\Gamma} T_{\alpha\beta}N^\alpha X^\beta\,dS
=\int_{D_t} T_{\alpha\beta}\nabla^\alpha X^\beta\,dV.
\]
Here $n$ is the future unit normal to the spacelike slices and $N$ is the outward normal to the timelike boundary. Because the slices are spacelike and $X$ is future timelike, the flux density
\[
T_{\alpha\beta}n^\alpha X^\beta
\]
is comparable to $|\partial u|^2$ on each slice. On the lateral boundary, the Dirichlet condition forces the tangential derivatives of $u$ along $\Gamma$ to vanish. Consequently the boundary flux has a good sign; more precisely,
\[
\int_{\Gamma} T_{\alpha\beta}N^\alpha X^\beta\,dS\ge 0
\]
for the outward choice of $N$.
Therefore
\[
E_X(t)\le E_X(0)+C\int_0^t E_X(s)\,ds,
\]
where
\[
E_X(t):=\int_{\Sigma_t} T_{\alpha\beta}n^\alpha X^\beta\,d\sigma.
\]
By Gronwall's inequality,
\[
E_X(t)\le e^{Ct}E_X(0),\qquad 0\le t\le T.
\]
Since $E_X(t)$ is equivalent to the natural $H^1\times L^2$ energy on each slice, we obtain the desired variable coefficient energy estimate. \prvd

\vspace{0.5em}

\subsection*{Problem 2}
Consider the semilinear wave equation in $\R^{3+1}$
\[
\Boxop u=u^3,\qquad u(0)=u_0,\qquad \partial_tu(0)=u_1.
\]
Prove that this problem is locally well-posed for initial data $(u_0,u_1)\in \dot H^1\times L^2$.

\vspace{0.5em}
\textbf{Sol.} \par
This is the energy-critical semilinear wave equation in dimension $3+1$, so the natural framework is the Strichartz theory for the wave equation. Let
\[
S(t)(u_0,u_1):=\cos(t|\nabla|)u_0+\frac{\sin(t|\nabla|)}{|\nabla|}u_1
\]
be the free wave evolution. Then a solution of
\[
\Boxop u=u^3
\]
should satisfy the Duhamel formula
\[
u(t)=S(t)(u_0,u_1)+\int_0^t \frac{\sin((t-s)|\nabla|)}{|\nabla|}\bigl(u(s)^3\bigr)\,ds.
\]
We use the standard wave Strichartz estimates. In dimension $3+1$, one has for the linear solution
\[
\|S(t)(u_0,u_1)\|_{L_t^3L_x^6(I\times\R^3)}+\|\nabla S(t)(u_0,u_1)\|_{L_t^\infty L_x^2(I\times\R^3)}+\|\partial_t S(t)(u_0,u_1)\|_{L_t^\infty L_x^2(I\times\R^3)}
\lesssim \|u_0\|_{\dot H^1}+\|u_1\|_{L^2}
\]
for every interval $I$. Also, if $F$ is a forcing term, then
\[
\left\|\int_0^t \frac{\sin((t-s)|\nabla|)}{|\nabla|}F(s)\,ds\right\|_{L_t^3L_x^6}
+\left\|\nabla\int_0^t \frac{\sin((t-s)|\nabla|)}{|\nabla|}F(s)\,ds\right\|_{L_t^\infty L_x^2}
\lesssim \|F\|_{L_t^1L_x^2}.
\]
Therefore the key is to control $u^3$ in $L_t^1L_x^2$. But
\[
\|u^3\|_{L_t^1L_x^2(I\times\R^3)}=\|u\|_{L_t^3L_x^6(I\times\R^3)}^3.
\]
So we define the solution space
\[
X_I:=C_t(I;\dot H^1)\cap C_t^1(I;L^2)\cap L_t^3(I;L_x^6).
\]
Fix $A>0$ and consider the ball
\[
B_A:=\left\{u\in X_I:\ \|u\|_{L_t^\infty\dot H^1}+\|\partial_tu\|_{L_t^\infty L^2}+\|u\|_{L_t^3L_x^6}\le A\right\}.
\]
Define the map
\[
\Phi(u)(t):=S(t)(u_0,u_1)+\int_0^t \frac{\sin((t-s)|\nabla|)}{|\nabla|}\bigl(u(s)^3\bigr)\,ds.
\]
By the above Strichartz estimates,
\[
\|\Phi(u)\|_{X_I}
\lesssim \|u_0\|_{\dot H^1}+\|u_1\|_{L^2}+\|u\|_{L_t^3L_x^6(I\times\R^3)}^3.
\]
To get a contraction we localize in time. Choose the interval $I=[0,T]$ so that the free evolution is small in the critical norm:
\[
\|S(t)(u_0,u_1)\|_{L_t^3L_x^6([0,T]\times\R^3)}\le \eta
\]
with $\eta>0$ to be chosen. This is possible because the linear solution belongs to $L_t^3L_x^6(\mathrm{loc})$ and the norm is absolutely continuous in time. Then if $u\in B_A$,
\[
\|\Phi(u)\|_{X_I}
\le C(\|u_0\|_{\dot H^1}+\|u_1\|_{L^2})+CA^3.
\]
Choosing $A\simeq \|u_0\|_{\dot H^1}+\|u_1\|_{L^2}$ and then taking $T$ small enough so that the $L_t^3L_x^6$ contribution is sufficiently small, we ensure that $\Phi$ maps $B_A$ to itself.

For the contraction estimate, let $u,v\in B_A$. Then
\[
\Phi(u)-\Phi(v)=\int_0^t \frac{\sin((t-s)|\nabla|)}{|\nabla|}\bigl(u^3-v^3\bigr)(s)\,ds,
\]
so
\[
\|\Phi(u)-\Phi(v)\|_{X_I}
\lesssim \|u^3-v^3\|_{L_t^1L_x^2}.
\]
Using
\[
|u^3-v^3|\le C(|u|^2+|u||v|+|v|^2)|u-v|
\]
and H\"older,
\[
\|u^3-v^3\|_{L_t^1L_x^2}
\lesssim \bigl(\|u\|_{L_t^3L_x^6}^2+\|v\|_{L_t^3L_x^6}^2\bigr)\|u-v\|_{L_t^3L_x^6}.
\]
If the interval is chosen so that $A^2$ times the implicit constant is $<1$, then $\Phi$ is a contraction on $B_A$. Hence by Banach's fixed-point theorem there exists a unique solution
\[
u\in C_t([0,T];\dot H^1)\cap C_t^1([0,T];L^2)\cap L_t^3([0,T];L_x^6).
\]
Continuous dependence on the data follows from the same contraction estimate applied to two solutions with nearby initial data. Thus the problem is locally well-posed in $\dot H^1\times L^2$. \prvd

\vspace{0.5em}

\subsection*{Problem 3}
Let $g_{\alpha\beta}$ be a pseudo-Riemannian metric in $\R^{n+1}$. Show that for any forward timelike vectors $X$ and $Y$ one has
\[
T_{\alpha\beta}X^\alpha Y^\beta\ge 0.
\]

\vspace{0.5em}
\textbf{Sol.} \par
Here $T_{\alpha\beta}$ denotes the stress-energy tensor of a scalar field:
\[
T_{\alpha\beta}=\partial_\alpha u\,\partial_\beta u-\frac12 g_{\alpha\beta}g^{\mu\nu}\partial_\mu u\,\partial_\nu u.
\]
We must prove the dominant energy inequality
\[
T_{\alpha\beta}X^\alpha Y^\beta\ge 0
\]
for future timelike $X,Y$.

Since this is a pointwise statement, we may work at one fixed point and choose local inertial coordinates there so that
\[
g_{\alpha\beta}=\mathrm{diag}(-1,1,\dots,1).
\]
By a Lorentz transformation we may further assume that
\[
X=(1,0,\dots,0).
\]
Write
\[
\partial u=(\partial_tu,\nabla u).
\]
Then
\[
T_{00}=\frac12\Bigl((\partial_tu)^2+|\nabla u|^2\Bigr),
\qquad
T_{0j}=\partial_tu\,\partial_ju.
\]
Now let $Y=(Y^0,Y')$ be any future timelike vector. Since $Y$ is future timelike, we have
\[
Y^0>0,
\qquad |Y'|<Y^0.
\]
Therefore
\[
T_{\alpha\beta}X^\alpha Y^\beta=T_{0\beta}Y^\beta=T_{00}Y^0+\sum_{j=1}^nT_{0j}Y^j.
\]
Using the formulas above,
\[
T_{\alpha\beta}X^\alpha Y^\beta
=\frac12\Bigl((\partial_tu)^2+|\nabla u|^2\Bigr)Y^0+\partial_tu\,\nabla u\cdot Y'.
\]
By Cauchy--Schwarz and the timelike property,
\[
\partial_tu\,\nabla u\cdot Y'\ge -|\partial_tu|\,|\nabla u|\,|Y'|
\ge -|\partial_tu|\,|\nabla u|\,Y^0.
\]
Hence
\[
T_{\alpha\beta}X^\alpha Y^\beta
\ge Y^0\left[\frac12\Bigl((\partial_tu)^2+|\nabla u|^2\Bigr)-|\partial_tu|\,|\nabla u|\right].
\]
But
\[
\frac12(a^2+b^2)-ab=\frac12(a-b)^2\ge 0.
\]
Applying this with $a=|\partial_tu|$ and $b=|\nabla u|$, we get
\[
T_{\alpha\beta}X^\alpha Y^\beta\ge 0.
\]
This proves the claim. \prvd

\vspace{0.5em}

\subsection*{Problem 4}
Consider the system
\[
\begin{cases}
 w_t+r_x=0,\\
 r_t-i(g+a)w=0,
\end{cases}
\]
where $a$ and $g$ are real constants, and the frequencies of $(w,r)$ are supported on the negative half-line. Equivalently,
\[
Pw=w,\qquad Pr=r,
\qquad P=\frac12(I-iH).
\]
Is the problem well-posed in the space of holomorphic functions $L^2\times \dot H^{1/2}$?

\vspace{0.5em}
\textbf{Sol.} \par
We denote
\[
\kappa:=g+a\in\R.
\]
Since the equations have constant coefficients, the cleanest way is to work in Fourier space. Let $\xi<0$ be a negative frequency. Then the system becomes
\[
\partial_t \widehat w + i\xi\widehat r=0,
\qquad
\partial_t \widehat r - i\kappa\widehat w=0.
\]
Equivalently,
\[
\partial_t
\begin{pmatrix}
\widehat w\\
\widehat r
\end{pmatrix}
=
A(\xi)
\begin{pmatrix}
\widehat w\\
\widehat r
\end{pmatrix},
\qquad
A(\xi)=
\begin{pmatrix}
0 & -i\xi\\
i\kappa & 0
\end{pmatrix}.
\]
Differentiating once more gives
\[
\partial_t^2\widehat w=\kappa\xi\,\widehat w,
\qquad
\partial_t^2\widehat r=\kappa\xi\,\widehat r.
\]
Because $\xi<0$, the sign of $\kappa\xi$ is opposite to the sign of $\kappa$.

\medskip
\emph{Case 1: $\kappa\ge 0$.} Then $\kappa\xi\le 0$ for all $\xi<0$, so every Fourier mode is oscillatory (or stationary if $\kappa=0$). In this regime the problem is well-posed.

Indeed, consider the energy
\[
E(t):=\kappa\|w(t)\|_{L^2}^2+\|r(t)\|_{\dot H^{1/2}}^2.
\]
When $\kappa>0$, this is equivalent to
\[
\|w(t)\|_{L^2}^2+\|r(t)\|_{\dot H^{1/2}}^2.
\]
We now show that it is conserved. First,
\[
\frac{d}{dt}\|r\|_{\dot H^{1/2}}^2
=2\Ree\int |\xi|\,\partial_t\widehat r(\xi)\,\ov{\widehat r(\xi)}\,d\xi.
\]
Using $r_t=i\kappa w$,
\[
\frac{d}{dt}\|r\|_{\dot H^{1/2}}^2
=2\kappa\Ree\int |\xi|\,i\widehat w(\xi)\,\ov{\widehat r(\xi)}\,d\xi.
\]
Since the support is in $\{\xi<0\}$, we have $|\xi|=-\xi$, hence
\[
\frac{d}{dt}\|r\|_{\dot H^{1/2}}^2
=-2\kappa\Ree\int \xi i\widehat w(\xi)\,\ov{\widehat r(\xi)}\,d\xi.
\]
On the other hand, using $w_t=-r_x$,
\[
\frac{d}{dt}\|w\|_{L^2}^2
=2\Ree\int \partial_tw\,\ov w\,dx
=-2\Ree\int r_x\,\ov w\,dx.
\]
Passing to Fourier variables,
\[
\frac{d}{dt}\|w\|_{L^2}^2
=-2\Ree\int i\xi\widehat r(\xi)\,\ov{\widehat w(\xi)}\,d\xi.
\]
Therefore
\[
\kappa\frac{d}{dt}\|w\|_{L^2}^2
=-2\kappa\Ree\int i\xi\widehat r(\xi)\,\ov{\widehat w(\xi)}\,d\xi.
\]
Taking complex conjugates inside the real part shows that this exactly cancels the derivative of $\|r\|_{\dot H^{1/2}}^2$. Hence
\[
\frac{d}{dt}E(t)=0.
\]
So for $\kappa>0$ we obtain
\[
\|w(t)\|_{L^2}+\|r(t)\|_{\dot H^{1/2}}
\lesssim
\|w(0)\|_{L^2}+\|r(0)\|_{\dot H^{1/2}},
\]
and standard semigroup theory gives existence and uniqueness.

If $\kappa=0$, then $r_t=0$ and $w_t=-r_x$, so
\[
r(t)=r(0),
\qquad
w(t)=w(0)-t\,r_x(0).
\]
Since $r\in \dot H^{1/2}$, we only know $r_x\in \dot H^{-1/2}$ in general, so the map fails to preserve $L^2$ unless one imposes extra regularity. Thus in the stated space $L^2\times \dot H^{1/2}$ the problem is not well-posed when $\kappa=0$.

\medskip
\emph{Case 2: $\kappa<0$.} Then for every $\xi<0$ we have $\kappa\xi>0$, so each Fourier mode solves
\[
\partial_t^2\widehat w=(\kappa\xi)\widehat w
\]
with positive coefficient. Hence the solutions contain exponentially growing factors
\[
e^{t\sqrt{\kappa\xi}}.
\]
Since $\sqrt{\kappa\xi}=\sqrt{|\kappa|\,|\xi|}$ grows like $|\xi|^{1/2}$ as $|\xi|\to\infty$, the high frequencies are violently unstable. In particular, one can choose initial data bounded in $L^2\times \dot H^{1/2}$ for which the solution norm becomes arbitrarily large in arbitrarily short time. Therefore the flow cannot be continuous in the stated topology, so the Cauchy problem is ill-posed.

To summarize: the problem is well-posed in $L^2\times \dot H^{1/2}$ only in the strictly positive case
\[
g+a>0.
\]
If $g+a\le 0$, then it is not well-posed in the stated space. \prvd

\end{document}
